% RESEARCH_DRAFT_SHA256: b9c9166c4ce9df6258e1148f1039c4ce64aab4497bd425c8bc97da6ed1573539
% Layout and source-reference wrappers expanded for the web renderer.
\documentclass[11pt]{article}
\usepackage[margin=25mm]{geometry}
\usepackage{fontspec,unicode-math}
\setmainfont{Latin Modern Roman}
\setmathfont{Latin Modern Math}
\usepackage{amsmath,amssymb,mathtools}
\usepackage{xurl,hyperref}
\hypersetup{colorlinks=true,urlcolor=blue,linkcolor=blue}
\setcounter{secnumdepth}{0}
\setlength{\parindent}{0pt}
\setlength{\parskip}{5pt}
\setlength{\emergencystretch}{3em}
\newcommand{\sref}[2]{\hyperlink{source:#1:#2}{[S#2]}}
\newcommand{\cataloguescope}[1]{\paragraph{Source question.}#1}
\newcommand{\floor}[1]{\left\lfloor#1\right\rfloor}
\begin{document}
\section*{Research-attempt provenance}
This standalone research attempt retains the complete original definition body below,
including its historical source and preparation notes. Those notes describe the
earlier catalogue preparation; the new attempt and its verification are identified
separately after that body. No original ID, catalogue status, or source statement is changed.

% UnsolvedMath ID 30006383; source catalogue code OWR-14299517-002
\section{Near-Linear List Decoding of AEL and Ta-Shma Codes}\label{30006383}
{\small UnsolvedMath ID 30006383 · Computer Science}
\subsection{Definitions and mathematical statement}

\paragraph{Shared notation.}Standard number systems are denoted by $\mathbb N,\mathbb Z,\mathbb Q,\mathbb R,\mathbb C$; index ranges, computational models and specialized notation follow the cited problem source and the definitions below.


List decoding returns all codewords within the allowed distance. Capacity refers to the source\textquotesingle s optimal rate-versus-error tradeoff; the Johnson bound is the standard distance-derived list-decoding radius. Near-linear time allows the stated logarithmic factors in block length.

A code $C\subseteq\Sigma^N$ is a set of permitted length-$N$ words. Its Hamming minimum distance is the minimum number of unequal coordinates of two distinct words. For a linear $[N,K,D]_q$ code, rate is $K/N$ and relative distance is $D/N$; decoding recovers a codeword from corrupted observations.

List decoding returns every codeword within a prescribed distance of the received word, allowing several answers. Efficient list decoding requires polynomial runtime and a controlled output-list size; the intended capacity and alphabet are those in the source.
\paragraph{Statement recorded in the supplied source.}
Can Alon–Edmonds–Luby codes be list-decoded up to capacity in linear time? Can Ta-Shma codes be list-decoded up to a radius approaching one half, or up to the Johnson bound, in near-linear time? \href{https://doi.org/10.4171/owr/2025/41}{[S2]}
\paragraph{Scope qualification.}
The rates, alphabets, capacity conventions, and dependence on gap-to-capacity are source-specific; exact linear and near-linear runtimes are separate requirements.
\paragraph{Source status.}
The supplied catalogue labels this record open. Its source-status metadata was checked against the supplied snapshot on 2026-10-11. This import does not independently establish current open status or mathematical correctness. The catalogue literature assessment is dated 2026-08-22.
\subsection{Short English statement}
Meet the stated optimal decoding radii in linear or near-linear time for both named code families.
\subsection{Sources}
\begin{itemize}
\item[S1] Ulam AI and the UnsolvedMath Contributors, UnsolvedMath: A Curated Collection of Open Mathematics Problems, record 30006383 (OWR-14299517-002). Catalogue curation and metadata: CC BY 4.0. The supplied source question is retained; definitions and the English summary are editorial additions.
\url{https://huggingface.co/datasets/ulamai/UnsolvedMath}.
\item[S2] New Mathematical Directions in Coding Theory (2026), Oberwolfach Reports, DOI 10.4171/owr/2025/41, pp. 2187--2242
\url{https://doi.org/10.4171/owr/2025/41}.
\item[S3] Srivastava--Tulsiani primary paper
\url{https://arxiv.org/abs/2504.20333}.
\item[S4] FOCS 2025 primary paper
\url{https://doi.org/10.1109/FOCS63196.2025.00077}.
\textit{Underlying Oberwolfach Reports material retains its source-specific licence, including CC BY-SA 4.0 where applicable; see THIRD\_PARTY\_NOTICES.md.}
\end{itemize}
\paragraph{Preparation and review.}
Original question attributed to the cited source and UnsolvedMath contributors. Codex (GPT-6; exact serving revision not exposed) assisted with definitions, English summaries, source inspection and formatting. No new human mathematical authorship or independent proof verification is asserted. Canonical ID 30006383 and source code OWR-14299517-002 checked against the complete supplied registry and existing repository statements; known flagged repetitions and obvious identical conjectures excluded. Full mathematical-equivalence/literature review has not been performed. Statement preparation and final review followed contribution prompts C and D; automated checks verify structure and identity, not open status.
\label{end:30006383}

\section{Literature and context check}
The primary report, printed pp. 2197--2199, already reports near-linear decoding results for AEL and some Ta-Shma radii, while asking about linear-time list decoding up to capacity and stronger radius/runtime combinations. These reported algorithms are context, not new results proved here. The following binary geometric bound and explicit stitching model investigate why a local-list argument alone is insufficient.
\section{Research attempt: List-size geometry and the global stitching bottleneck}

Let a binary code of length $n\geq1$ have relative distance at least $0<\delta\leq1/2$. Suppose $L$ distinct codewords lie within relative Hamming radius $0\leq\rho<1/2$ of a received word. Replace their bits by signs and divide by $\sqrt n$, obtaining unit vectors $v_1,\ldots,v_L$ and center $y$. Then
\[
\langle v_i,y\rangle\geq a:=1-2\rho,
\qquad \langle v_i,v_j\rangle\leq b:=1-2\delta\quad(i\ne j).
\]
Cauchy--Schwarz and direct expansion give
\[
L^2a^2\leq\Big\|\sum_i v_i\Big\|^2
\leq L+L(L-1)b.
\]
If $a^2>b$, division gives
\[
L\leq\frac{1-b}{a^2-b}=\frac{2\delta}{(1-2\rho)^2-(1-2\delta)}.
\]
The strict inequality is the Johnson-type radius restriction. At its boundary the denominator vanishes: this proof supplies no bound there or beyond it. Even inside this radius it only bounds how many codewords exist and does not locate them.

\paragraph{A complete restricted stitching algorithm.}
Suppose local lists $L_v$ of size at most $q$ are placed on the vertices of a finite tree, and each edge $uv$ specifies an explicitly testable allowed-pair relation $R_{uv}\subseteq L_u\times L_v$. Root the tree. At a leaf every label is feasible. At an internal vertex $v$, keep $a\in L_v$ precisely when every child $u$ has at least one retained label $b$ with $(a,b)\in R_{vu}$. Induction on subtree size proves that retained labels are exactly those extending to a consistent labeling of that subtree: the child extensions use disjoint subtrees and hence can be combined. Initializing the lists and computing all compatibility tests takes $O(|V|q+|E|q^2)$ time, including a one-vertex tree. Storing one witness per surviving parent-label/child pair reconstructs one global solution; enumerating all choices can of course take exponential time if exponentially many solutions exist.

The tree induction fails on a cycle because child regions can share constraints. For example, binary labels on a triangle with each edge requiring inequality have two choices at every vertex and every single-edge constraint is extendible, but no global assignment exists: the first two inequalities force the endpoints of the third edge equal. Conversely, if all edge relations allow every pair, local lists of size two give $2^{|V|}$ global combinations. Thus bounded local list size alone gives neither global feasibility nor a bounded number of candidates. The special code/expander structure must supply the missing restrictions. Exhaustively forming the Cartesian product of local lists does not produce a near-linear decoder.
\section{Verification and adversarial audit}
The geometric proof uses a positive center correlation and a strictly positive denominator; these conditions are stated. The dynamic program is proved only on trees with explicit pairwise constraints. The triangle example refutes this unrestricted local-consistency inference, not any property of the particular AEL construction. Outputting a list of $L$ length-$n$ words explicitly already needs $\Omega(nL)$ symbol writes.
\paragraph{Audit corrections.} Separated a list-cardinality theorem from an algorithm and kept the strict Johnson-radius denominator and tree restriction explicit.
\paragraph{External computation.} None; symbolic arguments and boundary checks only.
\section{Scope and remaining gaps}
No AEL or Ta-Shma decoder, stronger list bound, or improvement from near-linear to linear time is established. The precise unproved step is efficient global stitching using the target code's expansion beyond the elementary tree setting.
\paragraph{Integration review.} Codex (GPT-6; exact serving revision not exposed) reviewed the written arguments and expanded layout and source-reference wrappers. Repository metadata and the production web renderer were checked using Python 3.13 and JavaScript with MathJax 3.2.2. This model-assisted review is not independent human proof verification; no PDF compilation or formal proof checking is asserted.

\section{Completion Estimate}
\noindent\textbf{COMPLETION: 0\%}
The calculations below the preserved source statement investigate a concrete route,
but establish no substantiated new progress toward the entire original target.
Known special cases and failed sufficient conditions are not credited as a solution.
\section{References}
The original source references and catalogue attribution are preserved in the
Sources section above. The following URLs were actually consulted for this attempt;
retention of an additional original reference does not claim fresh consultation of it.
\begin{itemize}
\item \url{https://ems.press/content/serial-article-files/52245} (consulted 2026-10-11). 
\end{itemize}
\section{Model and contributor disclosure}
Codex GPT-6 completed written-proof review of this attempt at its declared scope; this is model-assisted review, not independent human proof checking.
OpenAI gpt-6-astra, requested reasoning configuration ultra, generated this research
attempt on 2026-10-11. The exact serving revision is not exposed. Codex GPT-6
performed orchestration and repository-format review; its exact serving revision
is also not exposed. The mathematical route, boundary checks and final adversarial
self-audit followed contribution prompts A then B in
\texttt{docs/contribution\_prompts.md}, after reading \texttt{CONTRIBUTING.md}.
Python 3.14.2 assembled the standalone documents; any mathematical computation is
specified in Verification. No theorem prover was used. No human mathematical
authorship, substantial human proof correction, or independent human proof checking
is asserted. Model-assisted review, source retrieval, and structural checks do not
establish correctness or novelty. Full private reasoning transcripts are unavailable
as reproducibility artifacts; the complete arguments, inputs and outcomes actually
used are stated here or in the companion batch audit.
\section{Final status}
Unresolved as a new research contribution. This is a concrete route investigation,
not a claim of new complete resolution or a current-status certification. A known
answer reproduced for context does not become a new contribution. The catalogue's
source-status label is preserved; this attempt status does not assert that a
source-labelled solved question is open.
\end{document}
